\documentclass[10.5pt,a4paper,twoside,twocolumn]{ctexart}

%============== 页面设置 ==============
\usepackage[top=2.54cm,bottom=2.54cm,left=2.2cm,right=2.2cm,
  headheight=13pt,headsep=3mm]{geometry}
\usepackage{graphicx}
\usepackage{subfig} 
\usepackage{url}
\usepackage{amsmath,amsthm}
\usepackage{amssymb,amsfonts}
\usepackage{booktabs}
\usepackage{float}
\usepackage{caption}
\usepackage{fancyhdr}
\usepackage{xcolor}
\usepackage{algorithm}
\usepackage{algpseudocode}
\usepackage[numbers]{natbib}
\usepackage[hidelinks]{hyperref}

%============== 字体 ==============
\setmainfont{Times New Roman}

%============== 页眉 ==============
\pagestyle{fancy}
\fancyhf{}
\fancyhead[CO]{\small 基于网格化时空聚合与大小模型协作的全生命周期检测方法}
\fancyhead[RO]{\small 顾作一 2515100292}
\fancyhead[LE]{\small 顾作一 2515100292}
\fancyhead[CE]{\small 基于网格化时空聚合与大小模型协作的全生命周期检测方法}
\renewcommand{\headrulewidth}{0.4pt}
\setlength{\headwidth}{\textwidth}

%============== 自定义命令 ==============
\newcommand{\T}{\mathsf{T}}
\newcommand{\R}{\mathbb{R}}
\newcommand{\N}{\mathbb{N}}
\newcommand{\bx}{\mathbf{x}}
\newcommand{\bc}{\mathbf{c}}
\newcommand{\bw}{\mathbf{w}}
\newcommand{\bh}{\mathbf{h}}
\newcommand{\bg}{\mathbf{g}}

%============== 定理环境 ==============
\newtheorem{definition}{定义}[section]
\newtheorem{theorem}{定理}[section]
\newtheorem{lemma}{引理}[section]

%============== 标题 ==============
\title{\heiti\zihao{2} 基于网格化时空聚合与大小模型协作的\\高速公路异常事件全生命周期检测方法}
\author{\fangsong\zihao{3} 顾作一 2515100292\\\zihao{6}（宁波大学，浙江 宁波 315000）}
\date{}

\begin{document}

\twocolumn[
\maketitle

\vspace{-6mm}
\noindent\rule{\textwidth}{0.4pt}

\vspace{2mm}
{\heiti\zihao{5} 摘\quad 要}\quad
\zihao{5}
针对传统跟踪式高速异常检测存在ID漂移重复报警、偶发检测失效误报、YOLO无法语义判别、事件消失难自动确认等痛点，本文提出网格化时空聚合方法。该方法采用YOLOv26s实时检测+网格引擎时空聚合+Qwen3.5 9B大模型异步确认的三级流水线架构，将画面划分为50×50像素网格，通过速度过滤与高斯加权累加构建空间热力图，引入梯度衰减与空间冷却机制，实现报警触发、异常确认、锁定与消警的全生命周期闭环。在三段真实监控视频（5696帧）上，无事件视频零误报，有事件视频Qwen调用较传统FSM减少90.5%和99.3%。

\vspace{3mm}
\noindent{\heiti\zihao{5} 关键词}\quad \zihao{5} 高速公路监控；异常事件检测；网格化时空聚合；视觉大模型；空间热力图


\vspace{3mm}
\begin{center}
\zihao{3}\textbf{Grid-Based Spatial-Temporal Aggregation with Small+Large Model for Full-Lifecycle Highway Anomaly Event Detection}\\
\vspace{2mm}
\zihao{5} GU Zuo-Yi 2515100292\\
\zihao{6} (Ningbo University, Ningbo 315000, China)
\end{center}

\zihao{5}\noindent\textbf{Abstract}\quad This paper proposes a grid-based spatial-temporal aggregation method to address ID-drift-induced repeated alarms, false triggers, limited semantic discrimination, and difficult auto-clearance in tracking-based highway anomaly detection. It employs a three-stage pipeline (YOLOv26s + grid engine + Qwen 3.5-9B VLM), partitioning frames into $50\times50$ grids and constructing a spatial heatmap via velocity filtering and Gaussian-weighted accumulation. Gradient decay and spatial cooldown enable full-lifecycle closed-loop management from alarm triggering to clearance. On three real surveillance videos (5,696 frames), the method achieves zero false alarms on event-free video and reduces Qwen calls by 90.5\% and 99.3\% vs traditional FSM.

\vspace{3mm}
\noindent\textbf{Keywords}\quad highway surveillance; anomaly event detection; grid-based spatial-temporal aggregation; vision-language model; spatial heatmap

\vfill

\noindent\rule{\textwidth}{0.4pt}
]

\clearpage

%====================================================================
\section{\heiti\zihao{4} 引言}

高速公路作为现代交通运输的主动脉，其运行安全直接关系到人民生命财产安全和经济社会运行效率\cite{parsa2019}。据统计，高速公路上约30\%的重大交通事故与异常停车（包括违停、抛锚、事故）直接相关\cite{wang2020}。据宁波市公路路网运行管理统计，全市高速公路月均异常事件报警量达到24.7万起，实际登记事件大约2万条，绝大多数都是重复报警和施工导致的报警，真实事件中违规停车和事故类事件占比较高\cite{highway2026}。因此，实现对高速公路异常停车事件的快速检测，重复事件的过滤以及事件类型的精细化识别具有重要的现实意义。

\subsection{行业痛点}

当前高速公路视频监控系统在实际部署中面临四大核心痛点：

\textbf{(1) 重复报警问题。} 基于个体目标跟踪（Tracker-based）的传统方法通过为目标分配唯一跟踪ID并维护状态机来实现异常判定\cite{wojke2017}。然而，大型车辆遮挡或检测器间歇性失效会导致跟踪器分配全新的ID，状态机的历史累积信息丢失，进而引发同一物理目标的反复报警。在本文实验中，传统FSM方法在15分钟视频片段中产生了26次Qwen调用，其中绝大多数指向同一违停车辆。

\textbf{(2) 偶发性检测失效引发的误报问题。} 高速公路场景中，正常行驶车辆在监控画面远端可能表现出极低的像素位移（每帧仅数像素），纯粹的位移阈值判定容易将其误判为静止目标。此外，环境光照变化、雨雾天气等因素会进一步加剧检测不稳定。

\textbf{(3) YOLO模型事件类型判别能力单一。} 通用目标检测模型（如YOLO系列）仅能输出目标类别（car、truck等）和边界框，无法提供语义级的事件理解。例如，社会车辆与施工车辆在外观上可能极为相似，但前者的路边停靠属于异常，后者的作业停靠属于正常，检测器本身无法做出这种区分。

\textbf{(4) 事件消失的自动确认难题。} 大多数现有方法聚焦于事件发生（onset）的检测，对事件结束（offset）的自动确认关注不足。当异常车辆驶离后，系统能否及时自动消警，直接关系到指挥预警信息的时效性和可信度。

\subsection{小模型+大模型的协作范式}

近年来，"小模型实时感知+大模型异步推理"的双层协作架构在智能监控领域展现出巨大潜力\cite{zheng2025,bai2023}。其核心思想是：由轻量级检测模型（如YOLO系列）承担高频实时目标感知任务，由大语言/视觉模型（如Qwen、GPT-4V）在检测到可疑事件后介入进行语义级判断。这种架构的关键挑战在于如何在保持检测灵敏度的前提下，最大限度降低大模型的调用频次——因为大模型单次推理通常需要数秒甚至数十秒，且消耗大量算力资源。

Cerberus\cite{zheng2025}通过运动掩码提示结合规则偏差检测实现了151倍加速，但其触发逻辑仍为静态规则。Mohanarangan等人\cite{mohanarangan2026}通过置信度感知的级联学习降低虚警率，但未引入时间维度的累积机制。上述方法的共同缺陷在于：大模型的调用决策缺乏充分的时间上下文信息，导致在"正常慢行"与"异常静止"之间无法做出可靠区分。

\subsection{本文贡献}

针对上述痛点，本文提出一种基于网格化时空聚合与大小模型协作的异常事件检测方法。其核心思想是将小模型的监控逻辑从"基于个体的追踪"彻底重构为"基于区域的时空聚合"——不再关注"谁"在停，只关注物理空间的占用状态是否异常，具体的事件分类由大模型解决。本文的主要贡献如下：

(1) 提出了网格化时空矩阵数据结构，将监控画面划分为$50\times50$像素的均匀网格，每个网格独立维护滞留计数和衰减状态；

(2) 设计了基于速度过滤与高斯加权累加的热力图构建算法，通过高斯平滑消除像素级检测抖动，利用速度阈值严格区分正常通行与潜在静止；

(3) 引入分阶段梯度衰减策略与空间冷却机制，实现从报警触发（count≥120）到消警确认（count<30）的完整闭环管理，有效消除重复报警；

(4) 集成YOLOv26s实时检测器与Qwen3.5 9B视觉大模型，构建"小模型前端感知—网格引擎时空聚合—大模型异步确认"的三级流水线架构，并将亮度自适应机制纳入检测前端以应对多样化光照条件。

%====================================================================
\section{\heiti\zihao{4} 相关工作}
%====================================================================

\subsection{高速公路异常事件检测}

高速公路异常事件检测长期以来是智能交通系统领域的研究热点。早期方法主要基于感应线圈和雷达等物理传感器\cite{chen2001}，虽然检测精度较高，但设备部署和维护成本高昂。随着视频监控的普及，基于计算机视觉的方法逐渐成为主流。Setchell等人\cite{setchell2001}提出了基于视觉的道路交通监控传感器，开创了视频检测的先河。近年来，深度学习方法在目标检测领域取得了突破性进展，YOLO系列\cite{redmon2016}和RT-DETR\cite{zhao2024}等模型在实时性和准确性之间取得了良好平衡。

在事件检测层面，Tahir等人\cite{tahir2023}提出了基于CCTV视频分析的实时事件驱动道路交通监控系统。Qu等人\cite{qu2024}利用多层级融合的时空特征进行高速公路事件检测。然而，这些方法大多聚焦于交通事故的瞬时检测，缺乏对事件全生命周期的管理能力。

\subsection{基于跟踪的异常检测}

ByteTrack\cite{zhang2022}是当前最主流的多目标跟踪算法之一，通过卡尔曼滤波预测和IoU匹配实现高效跟踪。基于ByteTrack的异常检测方法通常为每个跟踪ID维护独立的状态机\cite{liu2024}，通过累积目标在连续帧中的行为特征来判断是否异常。

然而，此类方法存在根本性缺陷。Zheng等人\cite{zheng2025}指出，在大目标遮挡场景下，ByteTrack的track\_buffer参数（默认30帧）无法覆盖长时遮挡，导致跟踪ID频繁切换。Mao等人\cite{mao2026}虽然利用贝叶斯推理引导VLM关注远场异常，但仍未解决ID继承问题。Xiao等人\cite{xiao2023}提出基于图结构的事件检测方法，通过目标间时空关系建模提高了检测鲁棒性，但图结构维护本身仍依赖稳定的跟踪ID。

\subsection{视觉大模型在监控中的应用}

随着视觉语言模型（VLM）的迅猛发展，其在视频异常检测中的应用日益广泛。Cerberus\cite{zheng2025}提出了级联式VLM框架，通过运动掩码提示实现实时视频异常检测。Mohanarangan等人\cite{mohanarangan2026}设计了置信度感知的级联学习框架用于弱监督异常检测。但上述方法均未引入空间维度的时间累积机制，导致正常过车与真正异常难以区分。

%====================================================================
\section{\heiti\zihao{4} 方法设计}
%====================================================================

\subsection{系统总体架构}

本文提出的系统由四个核心模块组成：前端感知模块、网格引擎模块、异步推理模块和可视化模块。系统架构如图\ref{fig:architecture}所示。

\begin{figure}[H]
\centering
\fbox{\begin{minipage}{0.85\columnwidth}
\vspace{2mm}
\begin{center}
\textbf{系统架构}\\
\vspace{1mm}
\begin{tabular}{c}
视频流(5fps) \\
$\downarrow$ \\
YOLOv26s + ByteTrack \\
$\downarrow$ \\
网格引擎(SpatialTemporalGrid) \\
$\downarrow$ \\
候选框聚类 $\rightarrow$ 异步Qwen \\
\end{tabular}
\end{center}
\vspace{2mm}
\end{minipage}}
\caption{系统总体架构}
\label{fig:architecture}
\end{figure}

\subsection{网格化时空矩阵}

\begin{definition}[网格化时空矩阵]
将分辨率为$W\times H$的视频帧划分为均匀的正方形网格，每个网格的边长为$S$像素。定义网格矩阵$\mathbf{G}\in\mathbb{N}^{R\times C}$，其中$R=\lceil H/S\rceil$，$C=\lceil W/S\rceil$，矩阵元素$g_{i,j}\in\mathbb{N}$表示网格$(i,j)$的当前滞留计数值。
\end{definition}

每个网格维护以下时空属性：
\begin{itemize}
\item $g_{i,j}^{\text{count}}\in\{0,1,\ldots,M\}$：滞留帧计数，须设置上限避免无限累加后无法及时消除；
\item $g_{i,j}^{\text{inactive}}\in\mathbb{N}$：连续无目标的帧数，用于梯度衰减。
\end{itemize}

\subsection{速度过滤与高斯加权累加}

\subsubsection{速度计算}

对于ByteTrack输出的每个检测目标$k$，设其在连续两帧中的中心点坐标分别为$(x_{t-1}^k,y_{t-1}^k)$和$(x_t^k,y_t^k)$，定义瞬时速度：

\begin{equation}
v_t^k = \sqrt{(x_t^k - x_{t-1}^k)^2 + (y_t^k - y_{t-1}^k)^2}
\end{equation}

\subsubsection{速度过滤规则}

考虑到摄像机监视距离250米，50~60km/h通过情况下，对于640*640检测窗口，5fps情况下，行进速度约为12像素/帧，保守设速度阈值为$V_{\text{th}}=8$像素/帧，。对每个目标$k$，检测框的宽度$w_k$和高度$h_k$（像素），计算其覆盖的网格范围半径$r_k = \lceil \max(w_k, h_k) / (2S) \rceil + 1$：

\begin{itemize}
\item 若$v_t^k > V_{\text{th}}$（正常通行），则对覆盖区域内所有网格执行衰减：
\begin{equation}
g_{i,j}^{\text{count}} \leftarrow \max(0, g_{i,j}^{\text{count}} - 3)
\end{equation}

\item 若$v_t^k \leq V_{\text{th}}$（潜在静止），则执行高斯加权累加。
\end{itemize}

\subsubsection{高斯加权累加}

首先目标中心网格累加值取4，取$g_{i,j}$上限为400，然后定义$3\times3$高斯核矩阵$\mathbf{K}$：

\begin{equation}
\mathbf{K} = \begin{bmatrix}
1 & 2 & 1 \\
2 & 4 & 2 \\
1 & 2 & 1
\end{bmatrix}
\end{equation}

对于目标中心所在的网格$(g_x,g_y)$及其$3\times3$邻域$(g_x+dx, g_y+dy)$，其中$dx,dy\in\{-1,0,1\}$，执行累加：

\begin{equation}
g_{g_y+dy,\,g_x+dx}^{\text{count}} \leftarrow \min\left(M,\; g_{g_y+dy,\,g_x+dx}^{\text{count}} + \mathbf{K}_{dy+1,\,dx+1}\right)
\end{equation}

此外，按目标尺寸扩展覆盖区域$r_k$，对邻域外的扩展网格施加距离加权：

\begin{equation}
g_{i,j}^{\text{count}} \leftarrow \min\left(M,\; g_{i,j}^{\text{count}} + \left\lfloor\frac{1}{1+d}\right\rfloor\right)
\end{equation}

其中$d = \sqrt{(i-g_y)^2 + (j-g_x)^2}$为网格中心到目标中心的欧氏距离。

\subsection{梯度衰减策略}

对于当前帧无静止目标的非活跃网格，为保证在30秒内检测到事件消失，且避免数据抖动造成qwen的大量触发，采用分阶段梯度衰减策略。设网格的连续无目标帧数为$f_{\text{inactive}}$：

\begin{equation}
\Delta(f) = \begin{cases}
1, & f \leq 50 \quad\text{(0--10秒)}\\
2, & 50 < f \leq 75 \quad\text{(10--15秒)}\\
4, & 75 < f \leq 100 \quad\text{(15--20秒)}\\
10, & f > 100 \quad\text{(20秒以上)}
\end{cases}
\end{equation}

每帧衰减：
\begin{equation}
g_{i,j}^{\text{count}} \leftarrow \max(0,\; g_{i,j}^{\text{count}} - \Delta(f_{\text{inactive}}))
\end{equation}

\subsection{空间聚类与候选框生成}

当某个网格的滞留计数达到报警阈值$T_{\text{hold}}=120$时即30帧6秒即可报警，加上大模型判断事件满足行业8-10秒的要求，该网格被标记为"热点"。收集所有热点网格的位置，通过形态学膨胀和连通域分析进行空间聚类：

\begin{algorithm}[H]
\caption{热点空间聚类算法}
\begin{algorithmic}[1]
\State \textbf{输入}: 热点网格集合 $\mathcal{H}=\{(x_i,y_i)\}$，当前时间 $t$
\State \textbf{输出}: 候选框列表 $\mathcal{C}$
\State 创建二值掩码 $\mathbf{M}\in\{0,1\}^{R\times C}$，$\forall (x,y)\in\mathcal{H}:\mathbf{M}_{y,x}\leftarrow 1$
\State $\mathbf{M}\leftarrow\text{Dilate}(\mathbf{M}, 3\times3\text{ kernel})$
\State $\text{contours}\leftarrow\text{FindContours}(\mathbf{M})$
\State $\mathcal{C}\leftarrow\emptyset$
\For{每个轮廓 $c\in\text{contours}$}
    \If{$\text{Area}(c) \geq 2$}
        \State $(x,y,cw,ch)\leftarrow\text{BoundingRect}(c)$
        \State $\mathcal{C}\leftarrow\mathcal{C}\cup\{(xS+cwS/2,\;yS+chS/2,\;cwS,\;chS)\}$
    \EndIf
\EndFor
\State \Return $\mathcal{C}$
\end{algorithmic}
\end{algorithm}

\subsection{闭环管理机制}

\begin{definition}[报警触发]
当候选框首次出现且不在空间冷却期内时，触发EVENT\_CONFIRM事件，调用Qwen进行异常确认。系统向Qwen发送双帧拼接图(图\ref{fig:EVENT})：左侧为T0时刻（20秒前）的锁定区域裁剪图，右侧为T20时刻（当前）的同一区域裁剪图。
\end{definition}

\begin{figure}[H]
	\centering
	\includegraphics[width=\columnwidth]{./1图1.jpg}
	\caption{事件确认双帧拼接图}
	\label{fig:EVENT}
\end{figure}

\begin{definition}[消警触发]
当候选框对应区域的网格滞留计数降至消警阈值$T_{\text{clear}}=30$以下时，触发RISK\_CLEARANCE事件，向Qwen发送锁定区域图片，询问区域内是否仍有车辆存在（后期可扩展至人员、火灾等类型），当qwen回答为false，视为消警成功。
\end{definition}

\begin{definition}[重触发拦截]
为防止计数在阈值附近振荡导致重复触发，引入require\_drop\_below\_30标志。首次EVENT\_CONFIRM触发后且qwen回答为true，该标志置为True；仅当计数降至30以且完成消警后，标志才重置为False，允许下一次触发。
\end{definition}

\subsection{空间冷却机制}

当Qwen判定某候选框为非异常事件时，将候选框中心坐标$(c_x,c_y)$加入全局冷却池，冷却半径为50像素，冷却时长为30秒。在冷却期内，任何新候选框若中心点落入已有冷却区域内，将被拦截而不触发报警，避免短期重复调用。

\subsection{亮度自适应检测}

高速公路监控场景的光照条件变化剧烈，从正午强光到夜间暗光，固定置信度阈值难以适应。本系统引入亮度自适应机制：计算每帧灰度均值$\bar{b}$，根据亮度梯级动态调整YOLO检测置信度阈值：

\begin{equation}
p_{\text{conf}}(\bar{b}) = \begin{cases}
0.08, & \bar{b} < 40 \quad\text{(极暗)}\\
0.1, & 40 \leq \bar{b} < 80 \quad\text{(昏暗)}\\
0.3, & \bar{b} \geq 80 \quad\text{(明亮)}
\end{cases}
\end{equation}

%====================================================================
\section{\heiti\zihao{4} 实验与分析}
%====================================================================

\subsection{实验设置}

实验环境：处理器为Intel Core i7，无GPU加速（CPU模式），内存16GB。检测模型为YOLOv26s，跟踪器为ByteTrack（track\_buffer=150, match\_thresh=0.9），视觉大模型为Qwen 3.5-9B（通过LM Studio本地部署），抽帧频率为5fps。

测试视频为真实高速公路监控录像（分辨率1920×1080，25fps原始帧率，h264编码），分别是test白天、能见度良好、无异常事件，test1夜间、能见度良好、有靠边停车，test2白天、能见度良好、有靠边停车。对比方法包括：(A) 基于ByteTrack跟踪ID的状态机方法（w/ FSM），(B) 基于ByteTrack跟踪的状态机方法+空间冷却，(C) 本文提出的网格化时空聚合方法。

\subsection{实验结果}

\begin{table}[H]
\caption{三种方法的Qwen调用次数对比}
\label{tab:qwen_calls}
\centering
\begin{tabular}{lcccc}
\toprule
视频 & FSM(ByteTrack) & +空间冷却 & 网格聚合(本文)\\
\midrule
test & 0 & 0 & \textbf{0}\\
test1 & 42 & 25 & \textbf{4}\\
test2 & 143 & 86 & \textbf{1}\\
\bottomrule
\end{tabular}
\end{table}

网格聚合方法在有事件的视频中共触发5次报警，其中3次是消失报警，相比FSM方法分别减少97.3\%和95.5\%的无效调用，有效消除了ID漂移和空间震荡导致的重复报警问题。

\begin{table}[H]
\caption{网格聚合方法全程实验结果}
\label{tab:full_results}
\centering
\begin{tabular}{lcccc}
\toprule
视频 & 实验帧数 & Qwen调用 & 报警次数 & 消警\\
\midrule
test & 3001 & 0 & 0 & --\\
test1 & 1938 & 4 & 1(违停) & 车已走\\
test2 & 757  & 1 & 1(违停) & 未触发\\
\bottomrule
\end{tabular}
\end{table}

上述实验数据与真实视频情况温和，基本验证了本文方法在真实场景中的有效性。

\subsection{热力图曲线}

图\ref{fig:test_curve}至图\ref{fig:test2_curve}展示了三个视频在网格聚合方法下的最大计数值变化曲线及Qwen调用时间点。

\begin{figure}[H]
\centering
\includegraphics[width=\columnwidth]{../results/test_curve.png}
\caption{test网格计数值变化曲线（亮光/无事件，最大值12）}
\label{fig:test_curve}
\end{figure}

\begin{figure}[H]
\centering
\includegraphics[width=\columnwidth]{../results/test1_qwen.png}
\caption{test1网格计数值变化曲线（暗光/有事件，含Qwen回复标注）}
\label{fig:test1_curve}
\end{figure}

\begin{figure}[H]
\centering
\includegraphics[width=\columnwidth]{../results/test2_curve.png}
\caption{test2网格计数值变化曲线（亮光/有事件，1次报警）}
\label{fig:test2_curve}
\end{figure}

\subsection{消融实验}

为验证各模块的独立贡献，设计以下消融变体：

\begin{table}[H]
\caption{消融实验（test1.mp4, 100帧）}
\label{tab:ablation}
\centering
\begin{tabular}{lcc}
\toprule
配置变体 & Qwen调用 & 相对变化\\
\midrule
完整系统 & 4 & 基准\\
去掉梯度衰减 & 12 & +200\%\\
去掉高斯加权 & 18 & +350\%\\
去掉空间冷却 & 7 & +75\%\\
去掉锁定冷却 & 16 & +300\%\\
\bottomrule
\end{tabular}
\end{table}

消融实验表明，梯度衰减、高斯加权和重触发拦截是三个最关键的模块，单独移除任何一个都会导致Qwen调用次数显著增加。

\subsection{讨论}

本文方法的核心优势在于将异常检测从"跟踪个体"转变为"聚合空间"。传统FSM方法的根本问题在于对跟踪ID的强依赖——当ByteTrack因遮挡重新分配ID时，状态机的历史累积信息全部丢失，导致同一物理位置的车辆被当作新目标重新计数，一般可以通过提高帧率解决，但提高帧率会导致算力需求增加，提升检测成本。而网格化方法直接操作物理空间的像素级网格，天然免疫于ID漂移，且受帧率影响小，抽帧频次下降只要调整相应阈值即可。

空间冷却机制通过在物理空间建立全局黑名单，进一步解决了因坐标微小漂移导致的同一位置重复触发问题。与FSM中将冷却绑定在track\_id上不同，空间冷却绑定在物理坐标上，不受跟踪ID变化的影响。

在"小模型+大模型"的架构层面，本文方法的独特之处在于：网格引擎充当时空信息"蓄水池"，在将决策权交给大模型之前，已经通过长时间的物理空间观察积累了充分的统计证据。这使得大模型的每次调用都建立在可靠的基础上，从而大幅降低了调用频次。

而且本项目直接采用了YOLO26s公共模型，未针对高速公路场景或夜间场景进行训练或微调，若做进一步微调将有助于提升夜间检测准确率，将进一步降低test1中的Qwen的调用频次。

\subsection{向缓行与拥堵识别的拓展应用}

本文提出的网格化时空聚合框架具有良好的场景扩展性。除异常停车检测外，该框架可自然地推广至以下应用场景：

\textbf{(1) 缓行检测。} 在拥堵或施工路段，车辆处于持续低速行驶状态，速度值稳定在$v\in[2,8]$像素/帧区间。此类目标的网格计数不会降至30以下（不触发消警），但也难以在短期内达到120（不触发报警），而是在中间区段形成持续稳定但低于阈值的"温热点"。通过监测网格计数的平均值和方差，可以区分"空荡路段"（count≈0）、"缓行区段"（30<count<120，方差小）和"违停区段"（count≥120，方差小）。

\textbf{(2) 拥堵识别。} 当大面积网格（如连续5×3及以上网格块）的计数同时处于50-100区间并维持超过60秒时，可判定为拥堵事件。此时系统可触发低优先级的拥堵预警，向Qwen发送广角全景图询问"该区域是否出现大面积车辆滞留"，而无需像违停事件那样进入高优先级的紧急响应流程。

\textbf{(3) 多事件分级管理。} 基于计数阈值和网格面积可以实现事件分级：count≥120且网格面积超过2-3辆车是触发一级报警（疑似事故），count≥120且网格面积为1辆车时触发二级预警（疑似违停），30<count<120且网格面积超过4-5辆车时触发三级预警（疑似缓行），count<30触发消警确认。这种分级机制使系统能够根据事件的严重程度差异化调度大模型资源，进一步优化算力分配。

%====================================================================
\section{\heiti\zihao{4} 结论}
%====================================================================

本文提出了一种基于网格化时空聚合的高速公路异常事件检测方法。该方法将监控画面划分为均匀网格，通过速度过滤和高斯加权累加构建空间热力图，利用梯度衰减和空间冷却实现闭环管理，并集成YOLOv26s检测器与Qwen视觉大模型构建异步检测流水线。实验结果表明，该方法相比传统基于跟踪ID的状态机方法，对于单一事件可极大减少重复报警次数，同时可减少对大模型调用次数，减少不必要的算力消耗。未来工作将探索将Grounding DINO等开放词汇检测模型集成到前端感知模块中，以进一步提升系统的泛化能力。

%====================================================================
% 参考文献
%====================================================================
\begin{thebibliography}{99}

\bibitem{wang2020}
Wang J, Wu J, Li Z, et al.
A survey on highway traffic incident detection.
\emph{IEEE Transactions on Intelligent Transportation Systems}, 2020, 21(11): 4489-4511.

\bibitem{redmon2016}
Redmon J, Divvala S, Girshick R, et al.
You only look once: Unified, real-time object detection.
\emph{Proceedings of the IEEE Conference on Computer Vision and Pattern Recognition (CVPR)}, 2016: 779-788.

\bibitem{wojke2017}
Wojke N, Bewley A, Paulus D.
Simple online and realtime tracking with a deep association metric.
\emph{Proceedings of the IEEE International Conference on Image Processing (ICIP)}, 2017: 3645-3649.

\bibitem{zhang2022}
Zhang Y, Sun P, Jiang Y, et al.
ByteTrack: Multi-object tracking by associating every detection box.
\emph{Proceedings of the European Conference on Computer Vision (ECCV)}, 2022: 1-21.

\bibitem{zhao2024}
Zhao Y, Lv W, Xu S, et al.
DETRs beat YOLOs on real-time object detection.
\emph{Proceedings of the IEEE/CVF Conference on Computer Vision and Pattern Recognition (CVPR)}, 2024: 16965-16974.

\bibitem{tahir2023}
Tahir M, Qiao Y, Kanwal N, et al.
Real-time event-driven road traffic monitoring system using CCTV video analytics.
\emph{IEEE Access}, 2023, 12: 181-194.

\bibitem{qu2024}
Qu Q, Shen Y, Yang M, et al.
Expressway traffic incident detection using a deep learning approach based on spatiotemporal features with multilevel fusion.
\emph{Journal of Transportation Engineering, Part A: Systems}, 2024, 150(3): 04024001.

\bibitem{liu2024}
Liu C, Chen J, Liu H, et al.
Toward efficient traffic incident detection via explicit edge-level incident modeling.
\emph{IEEE Internet of Things Journal}, 2024, 11(10): 18093-18106.

\bibitem{zheng2025}
Zheng Y, Shi X, Chen J, et al.
Cerberus: Real-time video anomaly detection via cascaded vision-language models.
\emph{arXiv preprint arXiv:2510.16290}, 2025.

\bibitem{mohanarangan2026}
Mohanarangan K, Palanisamy P.
Confidence-aware cascaded learning for real-time weakly supervised anomaly detection in surveillance video.
\emph{IEEE Access}, 2026, 14: 56388-56411.

\bibitem{mao2026}
Mao X, Sui B, Zhang W.
Zoom in, reason out: Efficient far-field anomaly detection in expressway surveillance videos via focused VLM reasoning guided by Bayesian inference.
\emph{arXiv preprint arXiv:2604.23724}, 2026.

\bibitem{bai2023}
Bai J, Bai S, Chu W, et al.
Qwen technical report.
\emph{arXiv preprint arXiv:2309.16609}, 2023.

\bibitem{parsa2019}
Parsa A B, Taghipour H, Derrible S, et al.
Real-time accident detection: Coping with imbalanced data.
\emph{Accident Analysis \& Prevention}, 2019, 129: 202-210.
DOI: 10.1016/j.aap.2019.05.014.

\bibitem{xiao2023}
Xiao D, Dianati M, Geiger W G, et al.
Review of graph-based hazardous event detection methods for autonomous driving systems.
\emph{IEEE Transactions on Intelligent Transportation Systems}, 2023, 24(5): 4697-4715.
DOI: 10.1109/TITS.2023.3240104.

\bibitem{setchell2001}
Setchell C, Dagless E L.
Vision-based road-traffic monitoring sensor.
\emph{IEE Proceedings-Vision, Image and Signal Processing}, 2001, 148(1): 78-84.

\bibitem{highway2026}
宁波市公路与运输管理中心.
关于公布2026年5月全市公路路网运行管理情况的通知, 附件4: 高速公路异常事件统计.
2026.

\bibitem{chen2001}
Chen C, Petty K, Skabardonis A, et al.
Freeway performance measurement system: Mining loop detector data.
\emph{Transportation Research Record}, 2001, 1748(1): 96-102.

\end{thebibliography}

\end{document}
